\exercisesection{\S~3}

\begin{enumerate}
\def\labelenumi{\arabic{enumi})}
\item
  设 G 是维数 \textgreater{} 0 的紧李群。则 G 的每个有限交换子群都是循环群，当且仅当 G 同构于 U、SU(2, C) 或 SU(2, C) 中某个极大环面的正规化子。
\item
  用 K 表示域 R、C 或 H 之一，用 n 表示一个整数 \(\geq 1\)。给空间 \(K_{s}^{n}\) 赋以通常的埃尔米特型。
\end{enumerate}

\begin{enumerate}
\def\labelenumi{\alph{enumi})}
\item
  证明 \(\mathbf{U}(n,\mathrm{K})\) 是紧实李群。
\item
  证明 \(\mathrm{K}_s^n\) 中半径为 1 的球面是 \(\mathbf{U}(n,\mathbf{K})\) 的（实）李齐性空间；一点的稳定子群同构于 \(\mathbf{U}(n - 1,\mathbf{K})\)。
\item
  推出群 \(\mathbf{U}(n,\mathbf{C})\) 与 \(\mathbf{U}(n,\mathbf{H})\) 都连通，而 \(\mathbf{O}(n,\mathbf{R})\) 有两个连通分支。
\item
  证明群 \(\mathbf{SO}(n,\mathbf{R})\) 与 \(\mathbf{SU}(n,\mathbf{C})\) 都连通。
\item
  证明群 \(\mathbf{O}(n,\mathbf{R})\)（相应地 \(\mathbf{U}(n,\mathbf{C})\)）是 Z/2Z（相应地 T）借助于 \(\mathbf{SO}(n,\mathbf{R})\)（相应地 \(\mathbf{SU}(n,\mathbf{C})\)）的半直积。
\end{enumerate}

\begin{enumerate}
\def\labelenumi{\arabic{enumi})}
\setcounter{enumi}{2}
\tightlist
\item
  \begin{enumerate}
  \def\labelenumii{\alph{enumii})}
  \tightlist
  \item
    证明实李群 \(\mathbf{U}(n, \mathbf{H})\) 的李代数是由诸如 \(x \in \mathrm{M}_n(\mathbf{H})\) 且满足 \(^t\bar{x} = -x\) 的矩阵组成的集合，括号为 \([x, y] = xy - yx\)。把它记作 \(\mathfrak{u}(n, \mathbf{H})\)。
  \end{enumerate}
\end{enumerate}

\begin{enumerate}
\def\labelenumi{\alph{enumi})}
\setcounter{enumi}{1}
\item
  把 \(\mathbf{C}\) 与子域 \(\mathbf{R}(i)\) 等同起来（它是 \(\mathbf{H}\) 的一个子域），并把 \(\mathbf{C}^{2n}\) 与 \(\mathbf{H}^n\) 借助同构 \((z_1, \ldots, z_{2n}) \mapsto (z_1 + jz_{n+1}, \ldots, z_n + jz_{2n})\) 等同起来。证明等式 \(\mathbf{U}(n, \mathbf{H}) = \mathbf{U}(2n, \mathbf{C}) \cap \mathbf{Sp}(2n, \mathbf{C})\)。
\item
  由 \(b)\) 推出，每个 \(\mathbf{C}_n\) 型的紧单（实）李代数都同构于 \(\mathfrak{u}(n, \mathbf{H})\)。
\end{enumerate}

\begin{enumerate}
\def\labelenumi{\arabic{enumi})}
\setcounter{enumi}{3}
\tightlist
\item
  设 \(n\) 是一个整数 \(\geq 1\)。
\end{enumerate}

\begin{enumerate}
\def\labelenumi{\alph{enumi})}
\item
  证明群 \(\mathbf{SU}(n,\mathbf{C})\) 是单连通的（利用习题 2）。
\item
  证明 \(\mathbf{SU}(n,\mathbf{C})\) 的中心由诸矩阵 \(\lambda .I_n\)（其中 \(\lambda \in \mathbf{C}\)，\(\lambda^n = 1\)）组成。
\item
  每个 \(A_{n}\) 型的概单紧李群都同构于 \(\mathbf{SU}(n+1,\mathbf{C})\) 关于由矩阵 \(\zeta^k. I_{n+1} (0 \leq k < d)\) 组成的循环子群的商群，其中 \(d\) 整除 \(n + 1\)，\(\zeta\) 是一个本原 \(d\) 次单位根。
\item
  证明群 \(\mathbf{SL}(n,\mathbf{C})\) 是单连通的（利用第 III 章，§6，节 9，定理 6）。
\end{enumerate}

\begin{enumerate}
\def\labelenumi{\arabic{enumi})}
\setcounter{enumi}{4}
\tightlist
\item
  设 \(n\) 为整数 \(\geq 1\)，用 \(\mathbf{Spin}(n, \mathbf{R})\) 表示对应于 \(\mathbf{R}^n\) 上通常二次型的既约 Clifford 群（《代数》第 IX 章，§9，节 5）。
\end{enumerate}

\begin{enumerate}
\def\labelenumi{\alph{enumi})}
\item
  证明 \(\mathbf{Spin}(n,\mathbf{R})\) 是紧李群，且满同态 \(\varphi:\mathbf{Spin}(n,\mathbf{R})\to\mathbf{SO}(n,\mathbf{R})\)（同上）是解析的，其核为 \(\{+1,-1\}\)。
\item
  证明当 \(n \geq 2\) 时，\(\mathbf{Spin}(n, \mathbf{R})\) 是连通且单连通的（利用习题 2）。群 \(\pi_{1}(\mathbf{SO}(n, \mathbf{R}))\) 是 2 阶循环群。
\item
  设 \(Z_{n}\) 是 \(\mathbf{Spin}(n,\mathbf{R})\) 的中心。证明当 n 为奇数时 \(Z_{n}=\{+1,-1\}\)，当 n 为偶数时 \(Z_{n}=\{+1,-1,\varepsilon,-\varepsilon\}\)，其中 \(\varepsilon=e_{1}\ldots e_{n}\) 是 \(R^{n}\) 的典则基中诸元素之积；当 r 为偶数（相应地奇数）时，群 \(Z_{2r}\) 同构于 \((\mathbf{Z}/2\mathbf{Z})^{2}\)（相应地 \(Z/4Z\)）。
\item
  证明每个 \(B_{n}\) 型（\(n \geq 2\)）的概单连通紧李群都同构于 \(\mathbf{Spin}(2n + 1, \mathbf{R})\) 或 \(\mathbf{SO}(2n + 1, \mathbf{R})\)。
\item
  若 r 为奇数（相应地偶数）且 \(\geq 2\)，则每个 \(D_{r}\) 型的概单连通紧李群同构于 \(\mathbf{Spin}(2r,\mathbf{R})\)、\(\mathbf{SO}(2r,\mathbf{R})\) 或 \(\mathbf{SO}(2r,\mathbf{R})/\{\pm I_{2r}\}\)（相应地同构于上述群之一或 \(\mathbf{Spin}(2r,\mathbf{R})/\{1,\varepsilon\}\)）。
\end{enumerate}

\begin{enumerate}
\def\labelenumi{\arabic{enumi})}
\setcounter{enumi}{5}
\tightlist
\item
  \begin{enumerate}
  \def\labelenumii{\alph{enumii})}
  \tightlist
  \item
    证明紧李群 \(\mathbf{U}(n,\mathbf{H})\) 连通且单连通（利用习题 2），且其中心为 \(\{\pm I_{n}\}\)。
  \end{enumerate}
\end{enumerate}

\begin{enumerate}
\def\labelenumi{\alph{enumi})}
\setcounter{enumi}{1}
\tightlist
\item
  每个 \(C_{n}\) 型（\(n \geq 3\)）的概单连通紧李群都同构于 \(\mathbf{U}(n, \mathbf{H})\) 或 \(\mathbf{U}(n, \mathbf{H}) / \{\pm I_{n}\}\)。
\end{enumerate}

\begin{enumerate}
\def\labelenumi{\arabic{enumi})}
\setcounter{enumi}{6}
\tightlist
\item
  设 A 是 Cayley 八元数代数（《代数》第 III 章，附录，节 3），并取同上中的基 \((e_i)_{0\leq i\leq 7}\)。用 V 表示由 \(e_1,\ldots ,e_7\) 生成的纯八元数子空间，用 E 表示 V 中由 \(e_1,e_2,e_3,e_5,e_6,e_7\) 生成的子空间。把 A 中由 \(e_0,e_4\) 生成的子代数与复数域 \(\mathbf{C}\) 等同起来，并用 G 表示幺元代数 A 的自同构拓扑群。
\end{enumerate}

\begin{enumerate}
\def\labelenumi{\alph{enumi})}
\item
  用 Q 表示由 Cayley 范数诱导的 V 上的二次型，使 \((e_{i})_{1\leq i\leq7}\) 是 V 的一个标准正交基。证明映射 \(\sigma\mapsto\sigma|V\) 是从 G 到群 \(\mathbf{SO}(\mathbf{Q})\) 的单同态，而该群同构于 \(\mathbf{SO}(7,\mathbf{R})\)。
\item
  证明 A 的乘法赋予 E 一个 3 维 C-向量空间结构，且 \(\{e_{1}, e_{2}, e_{3}\}\) 是它的一个基。用 \(\Phi\) 表示 E 上以该基为标准正交基的埃尔米特型。设 H 是 G 中 \(e_{4}\) 的稳定子群。证明映射 \(\sigma \mapsto \sigma|E\) 是从 H 到群 \(\mathbf{SU}(\Phi)\) 的同构，而该群同构于 \(\mathbf{SU}(3, \mathbf{C})\)。映射 \(\sigma \mapsto \sigma(e_{4})\) 诱导出 G/H 到 V 的球面中的一个嵌入，该球面同构于 \(S_{6}\)。
\item
  设 T 是 H 中由自同构 \(\sigma\)（满足 \(\sigma(e_{0}) = e_{0}\)，\(\sigma(e_{1}) = \alpha e_{1}\)，\(\sigma(e_{2}) = \beta e_{2}\)，\(\sigma(e_{3}) = \gamma e_{3}\)，\(\sigma(e_{4}) = e_{4}\)，\(\sigma(e_{5}) = \bar{\alpha} e_{5}\)，\(\sigma(e_{6}) = \bar{\beta} e_{6}\)，\(\sigma(e_{7}) = \bar{\gamma} e_{7}\)）组成的环面，其中 \(\alpha, \beta, \gamma\) 是三个模为 1 且满足 \(\alpha\beta\gamma = 1\) 的复数。设 N 是 T 在 G 中的正规化子；证明 N/T 的阶为 12（注意 N 的每个元素都必须使诸 \(\pm e_{i}, i \neq 0\) 组成的集合稳定）；由此推出 G 的秩为 2。
\item
  证明 G 是 \(G_{2}\) 型的连通半单群，且 G/H 可与 \(S_{6}\) 等同（证明 \(G_{0} \neq H\)；推出 \(G_{0}\) 是 \(G_{2}\) 型的，再用维数论证）。每个 \(G_{2}\) 型紧群都同构于 G。
\end{enumerate}

\begin{enumerate}
\def\labelenumi{\arabic{enumi})}
\setcounter{enumi}{7}
\item
  设 G 是 \(A_{n}, B_{n}, C_{n}, D_{n}\) 或 \(G_{2}\) 型的概单连通紧李群。证明 \(\pi_{2}(G) = 0\) 且 \(\pi_{3}(G) = Z\)（利用习题 2 至 7，以及 \(\pi_{i}(\mathbf{S}_{n})\) 当 i \textless{} n 时为零、当 i = n 时为循环群这一事实（参见~《一般拓扑学》第 XI 章）。
\item
  设 g 是紧李代数，t 是 g 的一个嘉当子代数；设 \((\mathrm{X}_{\alpha})_{\alpha\in\mathbb{R}}\) 是分裂约化代数 \((\mathfrak{g}_{\mathbf{C}},\mathfrak{t}_{\mathbf{C}})\) 中的一个谢瓦莱系，使得 \(X_{\alpha}\) 与 \(X_{-\alpha}\)（关于 g）对所有 \(\alpha\in R\) 共轭。
\end{enumerate}

\begin{enumerate}
\def\labelenumi{\alph{enumi})}
\item
  设 \(\mathcal{T}\) 是一个 \(\mathbf{Z}\)-子模，包含于 \(\mathfrak{t}_{\mathbf{C}}\)，它包含诸 \(iH_{\alpha}\)（\(\alpha \in \mathrm{R}\)），且满足 \(\alpha(\mathcal{T}) \subset \mathbf{Z}i\)（对所有 \(\alpha \in \mathrm{R}\)）。证明 \(\mathbf{Z}\)-子模 \(\mathcal{G}\)（\(\mathfrak{g}_{\mathbf{C}}\) 中由 \(\mathcal{T}\) 和诸元素 \(u_{\alpha}\)、\(v_{\alpha}\)（\(\alpha \in \mathrm{R}\)）生成的）是一个 \(\mathbf{Z}\)-李子代数（在 \(\mathfrak{g}\) 中）。
\item
  设 g（相应地 t）是紧群 G（相应地 G 的一个极大环面 T）的李代数。设 \(\Gamma(\mathrm{T})\) 是同态 \(\exp_{\mathrm{T}}:t\to\mathrm{T}\) 的核。证明 Z-模 \((2\pi)^{-1}\varGamma(\mathrm{T})\) 满足 a) 的假设。
\item
  设 \(\langle,\rangle\) 是 g 上的一个不变内积；设 \(\mu\)（相应地 \(\tau\)） 是当 g（相应地 t）借助一个标准正交基与空间 \(R^{n}\) 等同时，对应于 Lebesgue 测度的 g（相应地 t）上的 Haar 测度。再用 \(\mu\)（相应地 \(\tau\)） 表示 g/G（相应地 t/T）上由 \(\mu\)（相应地 \(\tau\)） 除以 G（相应地 T）上的规范化 Haar 测度所得的商测度。证明公式 \(\mu(\mathfrak{g}/\mathcal{G})=\tau(\mathfrak{t}/\mathcal{T}).\prod_{\alpha\in\mathrm{R}_{+}}\langle iH_{\alpha},iH_{\alpha}\rangle\)。
\end{enumerate}

